\documentclass[11pt]{article}

\usepackage[T1]{fontenc}
\usepackage{amsmath,amssymb,booktabs}
\usepackage{hyperref}

\title{Nested GLS Profile Spectrum: First Smoke Test}
\author{}
\date{\today}

\newcommand{\tr}{\operatorname{tr}}

\begin{document}
\maketitle

\section*{Purpose}

This note records what is needed to compute the full regular eigenspectrum for
nested tests in the profile-Hessian paradigm, and a first finite-sample smoke
test in a continuous GLS setting. The companion scripts are
\[
  \texttt{relative-and-nested-fit-check.py}
  \qquad\text{and}\qquad
  \texttt{gls\_nested\_profile\_mc.py}.
\]
The finite-sample script uses the C++ helper
\texttt{weighted\_chisq\_quantile\_cli.cpp}, which calls \texttt{magmaan}'s
\texttt{robust::weighted\_chisq\_quantile} and
\texttt{robust::weighted\_chisq\_upper} for deterministic mixture quantiles. The
only Monte Carlo in the smoke is the finite-sample data generation.

\section{Required Objects}

Let $x$ be the first-stage vector and let
\[
  \sqrt n(\hat x-x_0)\Rightarrow Z,\qquad Z\sim N(0,\Gamma_x).
\]
For each model $M$, define the profiled value
\[
  v_M(x)=\min_\theta f_M(\theta,x).
\]
The nested contrast is
\[
  c(x)=v_B(x)-v_A(x),
\]
where $B$ is the restricted model and $A$ is the larger model. Under the regular
nested pseudo-null,
\[
  c(x_0)=0,\qquad \nabla c(x_0)=0.
\]
The reference law is therefore
\[
  n\,c(\hat x)\Rightarrow \frac12 Z'(Q_B-Q_A)Z
\]
if $f$ is scaled as a half-discrepancy, or equivalently
\[
  2n\,c(\hat x)\Rightarrow Z'(Q_B-Q_A)Z.
\]
Thus the full spectrum is obtained from the self-adjoint version of
\[
  (Q_B-Q_A)\Gamma_x .
\]
In implementation terms, the steps are:
\begin{enumerate}
\item Build a common first-stage vector $x$ and covariance $\Gamma_x$.
\item Compute the profiled score $s_M(x)=\nabla_x v_M(x)$ for each model.
\item Differentiate it to get $Q_M=D_xs_M(x_0)$.
\item Verify the pseudo-null conditions $v_B=v_A$ and $s_B=s_A$.
\item Eigensolve $\Gamma_x^{1/2}(Q_B-Q_A)\Gamma_x^{1/2}$.
\end{enumerate}
The old $U_B-U_A$ spectrum is the special exact-fit/fixed-weight/Gauss--Newton
case. It is not the general object.

\section{GLS Smoke Design}

The smoke uses the linear covariance example from
\texttt{relative-and-nested-fit.tex}. There are four observed variables.
\[
  B:\ \text{compound symmetry: one variance and one covariance},
\]
\[
  A:\ \text{one common variance and free covariances}.
\]
The test is equality of the covariances. The normal-theory GLS weight is
estimated from the first-stage covariance:
\[
  W(u)=\Gamma_{\mathrm{NT}}(u)^{-1}.
\]

To make the nested pseudo-null exact while keeping both models misspecified, the
population covariance is diagonal with unequal variances:
\[
  \Sigma_0=I+\epsilon\operatorname{diag}(d),
  \qquad
  d=(1,-1,1,-1).
\]
Both models share the same pseudo-true zero covariance structure, so
\[
  v_B(x_0)-v_A(x_0)=0,\qquad s_B(x_0)-s_A(x_0)=0,
\]
but the common-variance part is misspecified when $\epsilon\ne0$.

\section{Population Spectrum}

For $\epsilon=0.15$, the population script reports:
\[
  v_B-v_A=0,\qquad \lVert s_B-s_A\rVert=0,
\]
and the profile/value finite-difference Hessians agree to $1.17\times10^{-7}$.
The classical fixed-weight spectrum is
\[
  (1,1,1,1,1),
\]
whereas the profile-Hessian spectrum is
\[
  (1.4082,\ 0.9139,\ 0.9139,\ 0.9139,\ 0.5057),
  \qquad
  \tr=4.6557 .
\]
So even in this linear model with no model-curvature term, estimating the GLS
weight changes the reference law under fixed misspecification.

\section{Finite-Sample Smoke}

The Monte Carlo uses $6000$ replications for each sample size, simulating normal
data from the same $\Sigma_0$ and refitting both GLS profiles with $W(\hat u)$.
The reference quantiles are computed once from the population spectrum.

\begin{center}
\begin{tabular}{lrrrr}
\toprule
Reference & q50 & q90 & q95 & q99 \\
\midrule
Profile spectrum & 3.9817 & 8.7316 & 10.5854 & 14.7734 \\
Classical $U$ / $\chi^2_5$ & 4.3515 & 9.2364 & 11.0705 & 15.0863 \\
\bottomrule
\end{tabular}
\end{center}

\begin{center}
\begin{tabular}{rrrrrr}
\toprule
$N$ & empirical q50 & empirical q90 & empirical q95 & tail at profile q95
  & tail at $\chi^2_5$ q95 \\
\midrule
150  & 3.8651 & 8.4670 & 10.1870 & 0.0437 & 0.0352 \\
300  & 3.9323 & 8.4929 & 10.3074 & 0.0457 & 0.0385 \\
600  & 3.9239 & 8.6244 & 10.4148 & 0.0475 & 0.0397 \\
1200 & 3.9793 & 8.7283 & 10.5150 & 0.0487 & 0.0405 \\
3000 & 3.8889 & 8.5904 & 10.5366 & 0.0487 & 0.0395 \\
\bottomrule
\end{tabular}
\end{center}

The profile cutoff is close to nominal once $N$ is moderate. The ordinary
$\chi^2_5$ cutoff remains conservative here because the profile spectrum has a
smaller upper tail cutoff than the exact-fit spectrum.

\section{Magmaan Machinery}

The reusable library pieces are now:
\[
  \texttt{robust::compute\_profile\_contrast\_spectrum}(Q,\Gamma),
\]
which returns the positive eigenvalues of
$\Gamma^{1/2}Q\Gamma^{1/2}$ and flags genuinely negative directions, and
\[
  \texttt{robust::weighted\_chisq\_upper},\qquad
  \texttt{robust::weighted\_chisq\_quantile},
\]
which provide the central positive weighted-\(\chi^2\) reference law and its
cutoffs. The old \texttt{robust::imhof\_upper} symbol remains as a compatibility
alias, but the new call sites use the distribution-level name.

\section{What This Suggests}

For GLS, the full nested eigenspectrum is not hard in principle. The hard part
is not the eigensolve; it is computing $Q_M$ reliably. Finite differences of the
profile score are enough for research scripts. The next production step is to
expose analytic or automatic profile derivatives for:
\[
  s_M(x),\qquad Q_M=D_xs_M(x).
\]

For categorical DWLS, the same steps apply, but $x$ must include the weight
inputs such as $\gamma=\operatorname{diag}(\Gamma_u)$. That makes $Q_M$ a block
Hessian over $(u,\gamma)$ rather than only over ordinal moments. The GLS smoke
is therefore the right first implementation target: it validates the profile
machinery before the ordinal first-stage stack is made larger.

\end{document}
